% Run pdfLaTeX from the source directory. Keep the lookup for the activity's year.
% Exactly ONE worked example, ONE guided example, and FIVE independent exercises.
% Counts apply to the whole activity; never repeat these blocks per family.
% Printable student activity only. Gray section/task headings, bold underlined
% subtask titles, white criterion boxes, and neutral guided response boxes are
% permitted structural/response-area styles. Semantic
% instructional highlights and legends are forbidden.
% Strip source highlighting while preserving content and ordinary formatting.
% Avoid inherited highlight-producing macros (e.g. the shaded preamble boxed macro).
\makeatletter
\def\input@path{{C:/Users/orlan/Desktop/SanBartolo/global_economics/book2026-2027/preamble/}}
\makeatother
\input{preamble.tex}
\usepackage{needspace}
\usepackage{varwidth}

\geometry{letterpaper,margin=0.25in}
\pagestyle{empty}
\linespread{1.04}
\setlength{\parskip}{4pt}
\raggedbottom

\begin{document}
\begingroup % Activity-only helpers; flush pages before ending this group.

% Exactly TWO table-style families. All helpers size to local \linewidth.
% Family 1: plain opening notation/formula reference, never a highlight legend.
% Adapt widths or use full width for the topic's actual definitions and formulas.
\newcommand{\NotationReferenceTable}[1]{%
    \par\begingroup
    \small
    \setlength{\tabcolsep}{3pt}
    \setlength{\extrarowheight}{0pt}
    \setlength{\parskip}{0pt}
    \renewcommand{\arraystretch}{1}
    \noindent\begin{tabularx}{\linewidth}{@{}>{\raggedright\arraybackslash}p{.21\linewidth}>{\raggedright\arraybackslash}p{.34\linewidth}>{\raggedright\arraybackslash}X@{}}
        \toprule
        \textbf{Symbol} & \textbf{Concept} & \textbf{Formula} \\
        \midrule
        #1
        \bottomrule
    \end{tabularx}\par
    \endgroup
}
% Family 2: all instructional/problem tables, in the printable activity.
% Width variants below are the SAME ruled family; adapt column count if needed.
% Body rows end with \\ \hline. Keep group/case headings short or wrap them.
\newcommand{\ClassroomRuledTable}[4]{%
    \par\begingroup
    \small
    \setlength{\tabcolsep}{4pt}
    \setlength{\extrarowheight}{0pt}
    \setlength{\parskip}{0pt}
    \renewcommand{\arraystretch}{1.16}
    \noindent\begin{tabularx}{\linewidth}{|>{\raggedright\arraybackslash}p{#1}|>{\raggedright\arraybackslash}p{#2}|>{\raggedright\arraybackslash}X|}
        \hline
        #3 \\ \hline
        #4
    \end{tabularx}\par
    \endgroup
}
\newcommand{\InformationTable}[1]{%
    \ClassroomRuledTable{.14\linewidth}{.32\linewidth}
        {\textbf{Symbol} & \textbf{Concept} & \textbf{In context...}}{#1}%
}
\NewDocumentCommand{\CalculationTable}{O{.10\linewidth} O{.20\linewidth} +m}{%
    \ClassroomRuledTable{#1}{#2}
        {\textbf{Symbol} & \textbf{Concept} & \textbf{Calculation, result, and meaning in context}}{#3}%
}
\newcommand{\TableGroup}[1]{\multicolumn{3}{|l|}{\textbf{#1}}\\ \hline}
% Structural styles for the printable activity only; no semantic highlighting.
\newcommand{\stageheading}[1]{%
    \par\Needspace{4\baselineskip}\addvspace{4pt}%
    {\centering\large\bfseries
     \tikz[baseline=(title.base)]{%
         \node[fill=black!10,draw=black,line width=0.2pt,text=black,rounded corners=2mm,
             inner sep=3pt,outer sep=0pt] (title) {\strut #1};}\par}%
    \nopagebreak[4]\smallskip
}
\newcommand{\criterionbox}[2]{%
    \begin{tcolorbox}[colback=white,colframe=black!45,boxrule=.4pt,
        arc=2mm,boxsep=0pt,left=6pt,right=6pt,top=5pt,bottom=5pt,
        before skip=4pt,after skip=4pt,width=\linewidth]
        \textbf{#1} #2
    \end{tcolorbox}%
}
\newcommand{\problemtitle}[3]{%
    \par\Needspace{4\baselineskip}\addvspace{0.4\baselineskip}%
    \noindent
    {\begin{varwidth}[t]{0.75\linewidth}%
         \bfseries\strut #1. #2%
     \end{varwidth}}%
    \hfill\textbf{#3}\par\nopagebreak[4]
}
% Level 2: full numbered task title in #1; optional instructions in #2.
% Shade only the title and padding; wrap within the local linewidth.
\newcommand{\taskheading}[1]{%
    \par\Needspace{5\baselineskip}\addvspace{4pt}%
    \noindent\tikz[baseline=(title.base)]{%
        \node[fill=black!5,draw=none,text=black,rounded corners=2mm,
            inner sep=3pt,outer sep=0pt] (title) {%
            \begin{varwidth}{\dimexpr\linewidth-6pt\relax}%
                \raggedright\bfseries\strut #1\strut
            \end{varwidth}};}\par\nopagebreak[4]\smallskip
}
\newcommand{\task}[2]{\taskheading{#1}#2}
% Level 3: restart 1--3 within each task; instructions stay outside underline.
% Keep titles short: ordinary underline is intentionally unbreakable.
\newcommand{\subtaskheading}[2]{%
    \par\Needspace{4\baselineskip}\addvspace{4pt}%
    \noindent{\color{black}\textbf{\underline{#1. #2.}}}\par
    \nopagebreak[4]\smallskip
}

{\centering
    {\Large\bfseries [Criterion and activity title]\par}
    Global Economics \textbar{} [Grade] \textbar{} [Term] \textbar{} [Lesson]\par
}
\criterionbox{[Criterion code]}{[Complete official criterion statement, verbatim.]}

\stageheading{Equations and basic procedure}
[Substantive plain formulas, notation, and conditions for the actual topic.]
% Every task and reusable procedure: at most THREE numbered steps.
\begin{enumerate}[label=\arabic*.,itemsep=2pt]
\item {[First consolidated procedure step.]}
\item {[Second consolidated procedure step.]}
\item {[Third consolidated procedure step, if needed.]}
\end{enumerate}\par
% Replace illustrative symbols/relations with actual topic content. Preserve all
% needed definitions and formulas; no unnecessary gaps between filled rows.
\noindent\begin{minipage}[t]{.49\linewidth}
\vspace{0pt}
\NotationReferenceTable{
    $a$ & First quantity & [Definition and units] \\
    $b$ & Second quantity & [Definition and units] \\
}
\end{minipage}\hfill
\begin{minipage}[t]{.49\linewidth}
\vspace{0pt}
\NotationReferenceTable{
    $T$ & Total & $T=a+b$ \\
    $d$ & Difference & $d=a-b$ \\
}
\end{minipage}\par

\stageheading{General tasks}
% Required mathematical outputs only; no commentary, justification, or checks.
% Derive actual outputs from the source; repeat for another task only if needed.
\task{Task 1 --- [Concise mathematical task title].}{}
\begin{enumerate}[label=\arabic*.,itemsep=2pt]
\item {[First group of required mathematical outputs.]}
\item {[Second group of required mathematical outputs.]}
\item {[Final outputs, including contextual interpretations where required.]}
\end{enumerate}
\task{Task 2 --- [Second mathematical task title, when needed].}{}
\begin{enumerate}[label=\arabic*.,itemsep=2pt]
\item {[Required calculations and contextual meanings.]}
\end{enumerate}

% Exactly one fully solved worked example, with clear step-by-step reasoning.
\stageheading{Worked example}
\problemtitle{1}{[Worked problem title]}{Type A}
[Context, units, and constraints. Illustrative givens: two quantities of 12 and
8 units. Replace this entire example with a complete source-aligned problem.]
\task{Task 1 --- Information tables.}{}
\subtaskheading{1}{Given and derived quantities}
\noindent\begin{minipage}[t]{.49\linewidth}
\vspace{0pt}
\InformationTable{
    \TableGroup{Given quantities}
    $a$ & First quantity & 12 units in the first part of the context. \\ \hline
    $b$ & Second quantity & 8 units in the second part of the context. \\ \hline
}
\end{minipage}\hfill
\begin{minipage}[t]{.49\linewidth}
\vspace{0pt}
\InformationTable{
    \TableGroup{Derived quantities}
    $T$ & Total & Combined quantity, in the same units. \\ \hline
    $d$ & Difference & Excess of the first quantity over the second, in units. \\ \hline
}
\end{minipage}\par
\task{Task 2 --- Totals, differences, and contextual meanings.}{}
\subtaskheading{1}{Calculations and meanings}
\CalculationTable{
    \TableGroup{Combined quantities}
    $T$ & Total & $T=a+b=12+8=20$ units.\par
      Together the two parts account for 20 units. \\ \hline
    $d$ & Difference & $d=a-b=12-8=4$ units.\par
      The first part exceeds the second by 4 units. \\ \hline
}
\subtaskheading{2}{Checks and interpretation}
$T-b=20-8=12=a$ and
$d+b=4+8=12=a$. Both recover the first quantity, and all values use the
same units. [Add every further step, check, and interpretation required by
the actual source problem; consolidation must not omit reasoning.]

% Exactly one guided example. Preserve its actual context, givens, values,
% task order, intermediate calculations, representations, explanations, and
% interpretations when updating; replace only response areas and layout.
\begingroup % Guided Example ONLY: helpers cannot reach other activity sections.
% Guided response boxes: neutral writing areas, never semantic highlights.
% Width and height are minimum dimensions; the invisible 0 anchors the baseline.
% Nothing visible belongs inside: no placeholders, underlines, hints, or answers.
\newcommand{\guidedresponsebox}[2][9mm]{%
    \tikz[baseline=(answer.base)]{%
        \node[fill=black!1,draw=black!25,line width=.3pt,
            rounded corners=2mm,inner sep=0pt,outer sep=.3pt,
            minimum width=#2,minimum height=#1] (answer) {\phantom{0}};}}
% Reserve measured prompt+pair height BEFORE its heading. Each call starts a
% fresh paired row; do not stack subsequent stages in unequal-height columns.
\NewDocumentCommand{\guidedpair}{+m +m}{%
    \par\noindent\begin{minipage}[t]{.49\linewidth}\vspace{0pt}\raggedright
        #1
    \end{minipage}\hfill
    \begin{minipage}[t]{.49\linewidth}\vspace{0pt}\raggedright
        #2
    \end{minipage}\par\addvspace{4pt}}
\stageheading{Guided Example}
\problemtitle{2}{[Guided problem title]}{Type A}
[Fresh context, givens, units, and constraints. Illustrative structure: two
cases share a first quantity of 15 units; second quantities are 6 and 9 units.
Replace with the actual topic's givens and case structure.]
\task{Task 1 --- Shared quantity, units, and meaning.}{}
\subtaskheading{1}{Shared information}
\InformationTable{
    \TableGroup{Shared information}
    $a$ & First quantity & 15 \guidedresponsebox{18mm}.\par
      Meaning: \par\nopagebreak[4]
      \guidedresponsebox[18mm]{\dimexpr\linewidth-1pt\relax} \\ \hline
}
% All dimensions are adjustable starting points, not universal handwriting sizes.
% Numbers: 12mm x 9mm; fractions: 16mm x 12mm; expressions: 35mm x 10mm;
% event names: 45mm x 10mm; interpretations: local width x 22mm or taller.
% Place all prompts, supplied steps, and units OUTSIDE the empty boxes.
\Needspace{22\baselineskip}
\task{Task 2 --- Totals, differences, and contextual meanings for each case.}{}
\guidedpair{
    \subtaskheading{1}{Case A calculations}
    \CalculationTable[.14\linewidth][.24\linewidth]{
        \TableGroup{Case A}
        $b$ & Second quantity & 6 units. [Contextual meaning.] \\ \hline
        $T$ & Total & $T=a+b=15+6=$\par\nopagebreak[4]
          \guidedresponsebox{12mm} \\ \hline
        $d$ & Difference & \guidedresponsebox[18mm]{\dimexpr\linewidth-1pt\relax} \\ \hline
    }
}{
    \subtaskheading{1}{Case B calculations}
    \CalculationTable[.14\linewidth][.24\linewidth]{
        \TableGroup{Case B}
        $b$ & Second quantity & 9 units. [Contextual meaning.] \\ \hline
        $T$ & Total & \guidedresponsebox[18mm]{\dimexpr\linewidth-1pt\relax} \\ \hline
        $d$ & Difference & \guidedresponsebox[18mm]{\dimexpr\linewidth-1pt\relax} \\ \hline
    }
}
% Fresh paired row aligns the next stage despite unequal calculation heights.
% Break at meaningful stage boundaries; never prescribe reference page numbers.
\Needspace{12\baselineskip}
\guidedpair{
    \subtaskheading{2}{Case A interpretation}
    [Required contextual meaning.]\par\nopagebreak[4]
    \noindent\guidedresponsebox[22mm]{\dimexpr\linewidth-1pt\relax}
}{
    \subtaskheading{2}{Case B interpretation}
    [Required contextual meaning.]\par\nopagebreak[4]
    \noindent\guidedresponsebox[22mm]{\dimexpr\linewidth-1pt\relax}
}
% When needed, add corresponding graphs in a fresh paired row within the same
% stage. Preserve every original representation and supplied intermediate step.

% Representative validation only: not part of the reusable scaffold.
\Needspace{30\baselineskip}
\task{Task 3 --- Response dimensions and corresponding representations.}{}
\subtaskheading{1}{Numerical and symbolic responses}
\guidedpair{
\textbf{Case A: calculations.}\par
Supply the count: $n=\guidedresponsebox{12mm}$ units.
\[
q=\guidedresponsebox[12mm]{16mm},\qquad
x=\guidedresponsebox[10mm]{35mm}
\]
Complete the numerator:
\[
q=\frac{\guidedresponsebox[9mm]{12mm}}{12}.
\]
}{
\textbf{Case B: calculations.}\par
Supply the count: $n=\guidedresponsebox{12mm}$ units.
\[
\begin{aligned}
q&=\guidedresponsebox[12mm]{16mm}\\[4pt]
x&=\guidedresponsebox[10mm]{35mm}\\[4pt]
u&=\frac{3}{\guidedresponsebox[9mm]{12mm}}
\end{aligned}
\]
}
\Needspace{22\baselineskip}
\subtaskheading{2}{Events and graphs}
\guidedpair{
\textbf{Case A: event and graph.}\par
Name the event: \guidedresponsebox[10mm]{45mm}.
\par\nopagebreak[4]
\begin{tikzpicture}[x=7mm,y=7mm]
\draw[step=1,black!15] (0,0) grid (6,3);
\draw[->] (0,0)--(6.5,0) node[right] {$t$};
\draw[->] (0,0)--(0,3.5) node[above] {$v$};
\end{tikzpicture}
}{
\textbf{Case B: event and graph.}\par
Name the event: \guidedresponsebox[10mm]{45mm}.
\par\nopagebreak[4]
\begin{tikzpicture}[x=7mm,y=7mm]
\draw[step=1,black!15] (0,0) grid (6,3);
\draw[->] (0,0)--(6.5,0) node[right] {$t$};
\draw[->] (0,0)--(0,3.5) node[above] {$v$};
\end{tikzpicture}
}
\Needspace{12\baselineskip}
\guidedpair{
\subtaskheading{3}{Case A meaning}
Interpret the quantity in context over repeated observations.
\par\nopagebreak[4]\noindent
\guidedresponsebox[25mm]{\dimexpr\linewidth-1pt\relax}
}{
\subtaskheading{3}{Case B meaning}
Interpret the quantity in context over repeated observations.
\par\nopagebreak[4]\noindent
\guidedresponsebox[25mm]{\dimexpr\linewidth-1pt\relax}
}

\par\clearpage % Flush Guided Example before restoring all helper definitions.
\endgroup
\ifdefined\guidedresponsebox\errmessage{Guided response helper leaked}\fi
\ifdefined\guidedpair\errmessage{Guided pair helper leaked}\fi

% Exactly five UNSOLVED independent exercises: descriptions and givens only.
% No repeated tasks, procedural hints, starter calculations, checklists, response
% lines, blank tables, graph templates, reserved workspace, or solving locations.
% Every description supplies context, period where relevant, definitions, values,
% units, constraints, assumptions, and any essential problem-specific output.
% A compact given-data table is permitted. Do not prescribe lesson-specific math.
% Problems 3--5 consolidate core skills; problems 6--7 add modest challenge.
\stageheading{Independent Practice}
\problemtitle{3}{[First independent problem title]}{Type A}
[Fresh self-contained problem description with all required givens.]

\problemtitle{4}{[Second independent problem title]}{Type A}
[Fresh self-contained problem description with all required givens.]

\problemtitle{5}{[Third independent problem title]}{[Family label]}
[Fresh self-contained description consolidating core skills, with all givens.]

\problemtitle{6}{[Fourth independent problem title]}{[Family label]}
[Fresh self-contained description with all givens; modest challenge within the topic and criterion.]

\problemtitle{7}{[Fifth independent problem title]}{[Family label]}
[Fresh self-contained description with all givens; modest challenge within the topic and criterion.]

\par\clearpage % Flush the printable activity before restoring other styles.
\endgroup
% Preserve any existing appended solutions here, with their original formatting.
% Test that appended material cannot see the response helpers.
\ifdefined\guidedresponsebox\errmessage{Response helper reached appended material}\fi
\ifdefined\guidedpair\errmessage{Pair helper reached appended material}\fi
\end{document}
